\section{Preliminaries}\label{sec:preliminaries}

This section fixes notation and collects the facts used in the proofs: basic
properties of graphons and of the two variational problems, the first
variation of homomorphism densities, the rooted densities that describe the
effect of inserting a new block of vertices, and two results from analysis in
Banach spaces.

\subsection{Graphs and graphons}\label{sec:graphons}

Throughout, $H$ is a finite simple graph with vertex set $V(H)$ and edge set
$E(H)$, and we write $N:=|V(H)|$ and $m:=|E(H)|\ge1$.  Isolated vertices do not
affect homomorphism densities, so we assume that $H$ has none.  For $v\in V(H)$,
we write $N_H(v)$ for the set of neighbors of $v$ and $d_v:=|N_H(v)|$ for its
degree, and $\Delta(H)$ for the maximum degree.  Thus $\Delta(H)=1$ exactly
when $H$ is a perfect matching.

A graphon is a symmetric measurable function $W:[0,1]^2\to[0,1]$; we identify
graphons that agree almost everywhere, and $\W_0$ denotes the space of
graphons.  Recall the edge density $e(W)$, the homomorphism density $t(H,W)$,
and the relative entropy $I_p(W)$ from~\eqref{eq:graphon-quantities}.  We use
the cut distance $\delta_\cut$ and the quotient space $\widetilde\W_0$ as
in~\cite{LovaszGraphLimits}: the space $(\widetilde\W_0,\delta_\cut)$ is
compact, $e$ and $t(H,\cdot)$ are continuous with respect to $\delta_\cut$,
and $I_p$ is lower semicontinuous~\cite{ChatterjeeVaradhanLDP}.  All objective
functions and constraints below are invariant under measure-preserving
relabelings of $[0,1]$.

Unless stated otherwise, equalities and inequalities between measurable
functions hold almost everywhere.  For functions $u_1,u_2$ on $[0,1]$, we write
$u_1\wedge u_2$ and $u_1\vee u_2$ for their pointwise minimum and maximum.  For a graphon $W$ and $x\in[0,1]$, the
function
\[
  W_x(y):=W(x,y)\qquad(y\in[0,1])
\]
is the \emph{row} of $W$ at $x$.  By Fubini's theorem, $W_x$ is measurable for
almost every $x$, and two rows are regarded as equal if they agree almost
everywhere.  A graphon is \emph{multipodal} if there are a finite measurable
partition $[0,1]=A_1\cup\dots\cup A_M$ and numbers $b_{ij}=b_{ji}\in[0,1]$ such
that $W=b_{ij}$ almost everywhere on $A_i\times A_j$ for all $i,j$; the sets
$A_i$ are its \emph{blocks}, or \emph{podes}.  Equivalently, $W$ is the
graphon of a stochastic block model with finitely many blocks.  A multipodal
graphon has finitely many distinct rows.  We say that the blocks of a
multipodal graphon are \emph{ordered} if, after relabeling the blocks,
\begin{equation}\label{eq:ordered-blocks}
  b_{ij}\le b_{i'j}\qquad\text{whenever }i\le i',\ \text{for every }j;
\end{equation}
by symmetry, the matrix $(b_{ij})$ is then nondecreasing in both indices.
Equivalently, the distinct rows $g_1,\dots,g_M$ of $W$ satisfy
$g_1\le g_2\le\dots\le g_M$ pointwise almost everywhere.

\subsection{Relative entropy and graphon entropy}

For $q\in(0,1)$ and $z\in[0,1]$, recall
$J_q(z)=z\log\frac zq+(1-z)\log\frac{1-z}{1-q}$, with $0\log0=0$.  We write
\[
  \lo(z):=\log\frac z{1-z},
  \qquad
  \sigma(\xi):=\frac1{1+e^{-\xi}}
\]
for the log-odds function and its inverse.  For $0<z<1$,
\[
  J_q'(z)=\lo(z)-\lo(q),
  \qquad
  J_q''(z)=\frac1{z(1-z)}\ge4.
\]
Thus $J_q'$ is strictly increasing from $-\infty$ to $+\infty$ on $(0,1)$ and
vanishes only at $z=q$.  Recall also the entropy density and the graphon
entropy
\[
  S_0(z)=-\frac12\bigl(z\log z+(1-z)\log(1-z)\bigr),
  \qquad
  s(W)=\int_{[0,1]^2}S_0\bigl(W(x,y)\bigr)\dd x\dd y.
\]
A direct computation gives $J_q(z)=-2S_0(z)-\log(1-q)+z\log\frac{1-q}q$, so
\begin{equation}\label{eq:entropy-identity}
  I_\eps(W)=2S_0(\eps)-2s(W)\qquad\text{whenever }e(W)=\eps.
\end{equation}
Hence, among graphons with a prescribed edge density $\eps$, minimizing
$I_\eps$ is equivalent to maximizing $s$.

\subsection{Derivatives of homomorphism densities}\label{sec:derivatives}

We regard $e$ and $t(H,\cdot)$ as polynomial functions on the Banach space
$L^\infty_{\mathrm{sym}}([0,1]^2)$ of bounded symmetric measurable functions.
For an edge $ab\in E(H)$, oriented from $a$ to $b$, and a graphon $W$, let
\[
  K_{ab,W}(x,y):=\int_{[0,1]^{V(H)\setminus\{a,b\}}}
  \prod_{ij\in E(H)\setminus\{ab\}}W(x_i,x_j)
  \prod_{i\in V(H)\setminus\{a,b\}}\dd x_i,
  \qquad x_a=x,\ x_b=y.
\]
Reversing the orientation transposes this function, so the kernel
\begin{equation}\label{eq:derivative-kernel}
  \mathcal D_W(x,y):=\frac12\sum_{ab\in E(H)}\bigl(K_{ab,W}(x,y)+K_{ab,W}(y,x)\bigr)
\end{equation}
is well defined and symmetric.  It represents the first variation of the
homomorphism density: for every bounded symmetric $U$,
\begin{equation}\label{eq:first-variation}
  Dt(H,W)[U]=\int_{[0,1]^2}\mathcal D_W(x,y)\,U(x,y)\dd x\dd y.
\end{equation}
Since each $K_{ab,W}$ takes values in $[0,1]$, we have $0\le\mathcal D_W\le m$.
If moreover $W\ge\omega>0$, then $\mathcal D_W\ge m\omega^{m-1}$.

\subsection{Rooted densities}\label{sec:rooted-densities}

The central object of this paper describes how the homomorphism density
changes when a new type of vertex is added to a graphon.  Let $W$ be a graphon
and let $u:[0,1]\to[0,1]$ be measurable.  Think of $u$ as the row of a new
vertex: $u(y)$ is the probability that the new vertex is adjacent to a vertex
with label $y$.  For $v\in V(H)$, the \emph{rooted density} of $H$ at $v$ is the
density of copies of $H$ in which the new vertex plays the role of $v$:
\begin{equation}\label{eq:rooted-density}
  \mathcal R_{v,W}(u):=\int_{[0,1]^{V(H)\setminus\{v\}}}
  \prod_{z\in N_H(v)}u(x_z)
  \prod_{\substack{ab\in E(H)\\ a,b\ne v}}W(x_a,x_b)
  \prod_{z\in V(H)\setminus\{v\}}\dd x_z,
\end{equation}
and we put
\begin{equation}\label{eq:total-rooted-density}
  \mathcal R_W(u):=\sum_{v\in V(H)}\mathcal R_{v,W}(u).
\end{equation}
If $u=W_x$ is a row of $W$, then $\mathcal R_{v,W}(W_x)$ is the homomorphism
density of $H$ in $W$ with the vertex $v$ fixed at $x$, a rooted homomorphism
density in the sense of~\cite{LovaszGraphLimits}; the definition above extends
it to rows that do not belong to $W$.  Integrating $\mathcal R_{v,W}(W_x)$ over
$x$ restores the missing variable, so
\begin{equation}\label{eq:rooted-average}
  \int_0^1\mathcal R_W(W_x)\dd x=N\,t(H,W).
\end{equation}
The map $u\mapsto\mathcal R_W(u)$ is a polynomial on $L^\infty([0,1])$ of
degree at most $\Delta(H)$, and each of its monomials has a nonnegative kernel.
Its first derivative is represented by a bounded function,
$D\mathcal R_W(u)[h]=\int_0^1\nabla\mathcal R_W(u)(y)h(y)\dd y$, where
$\nabla\mathcal R_W(u)$ is the sum over the ordered pairs $(v,z_0)$ with
$z_0\in N_H(v)$ of
\begin{equation}\label{eq:rooted-gradient}
  G_{v,z_0}(u)(y):=\int\prod_{z\in N_H(v)\setminus\{z_0\}}u(x_z)
  \prod_{\substack{ab\in E(H)\\ a,b\ne v}}W(x_a,x_b)
  \prod_{z\in V(H)\setminus\{v,z_0\}}\dd x_z,
  \qquad x_{z_0}=y.
\end{equation}
There are $2m$ such pairs and each $G_{v,z_0}(u)$ takes values in $[0,1]$, so
\begin{equation}\label{eq:rooted-lipschitz}
  0\le\nabla\mathcal R_W(u)\le2m,
  \qquad
  |\mathcal R_W(u)-\mathcal R_W(u')|\le2m\|u-u'\|_1
\end{equation}
for measurable $u,u'$ with values in $[0,1]$.  Substituting $u=W_x$ and pairing
the two orientations of each edge shows that
\begin{equation}\label{eq:rooted-gradient-derivative}
  \nabla\mathcal R_W(W_x)(y)=2\,\mathcal D_W(x,y)
  \qquad\text{for almost every }(x,y).
\end{equation}

\begin{expl}\label{expl:triangle-rooted}
For $H=K_3$, one has
$\mathcal R_W(u)=3\int_{[0,1]^2}W(y,z)u(y)u(z)\dd y\dd z$.  If $u_1$ and $u_2$ are
two rows and $h_1:=(u_1-u_2)_+$ and $h_2:=(u_2-u_1)_+$, then
\[
  \mathcal R_W(u_1\wedge u_2)+\mathcal R_W(u_1\vee u_2)
  -\mathcal R_W(u_1)-\mathcal R_W(u_2)
  =6\int_{[0,1]^2}W(y,z)\,h_1(y)h_2(z)\dd y\dd z.
\]
When $W$ is bounded below by a positive constant, the right-hand side is
positive whenever $u_1$ and $u_2$ cross, that is, whenever both
$\{u_1>u_2\}$ and $\{u_2>u_1\}$ have positive measure.  This strict
supermodularity is the mechanism behind the ordering of rows.
\end{expl}

\subsection{Two results from analysis}

We use the following two standard results.  The first is the Lagrange
multiplier rule in Banach spaces; see, for example,~\cite{Zeidler1995}.  Its
surjectivity hypothesis is the classical constraint qualification, and the
condition in Definition~\ref{def:constraint-qualification} below is a
strengthened form of it.

\begin{lemma}[Lagrange multipliers]\label{lem:lagrange}
Let $X$ be a Banach space, $\mathcal O\subseteq X$ open, and let
$f:\mathcal O\to\RR$ and $g:\mathcal O\to\RR^k$ be continuously Fr\'echet
differentiable.  If $x_0\in\mathcal O$ is a local minimizer of $f$ on
$\{x\in\mathcal O:g(x)=g(x_0)\}$ and $Dg(x_0):X\to\RR^k$ is surjective, then
there is $\lambda\in\RR^k$ with $Df(x_0)=\lambda\cdot Dg(x_0)$.
\end{lemma}

The second is the analytic implicit function theorem in Banach
spaces~\cite{Whittlesey1965}.

\begin{lemma}[Analytic implicit function theorem]\label{lem:analytic-ift}
Let $X,Y,Z$ be real Banach spaces, let $\mathcal O\subseteq X\times Y$ be
open, and let $\Xi:\mathcal O\to Z$ be real-analytic with
$\Xi(x_0,y_0)=0$.  If $D_y\Xi(x_0,y_0):Y\to Z$ is a bounded isomorphism,
then there are neighborhoods $\mathcal U\ni x_0$ and $\mathcal V\ni y_0$ and a
real-analytic map $\varphi:\mathcal U\to\mathcal V$ such that, for
$(x,y)\in\mathcal U\times\mathcal V$, $\Xi(x,y)=0$ if and only if
$y=\varphi(x)$.
\end{lemma}
